\section{Waves: wave kinetics and quasilinear theory}
\label{sec:waves}

\subsection{Wave turbulence}

For a weakly nonlinear dispersive system with four-wave interactions and
dispersion relation $\omega_k$, the random-phase closure gives the wave
kinetic equation for the spectrum $n_k$ \cite{ZLF92},
\begin{equation}
  \partial_tn_1\;\propto\;\lambda^2\!\int n_1n_2n_3n_4
  \Bigl(\frac1{n_1}+\frac1{n_2}-\frac1{n_3}-\frac1{n_4}\Bigr)
  \delta(k_1+k_2-k_3-k_4)\,\delta(\omega_1+\omega_2-\omega_3-\omega_4)
  \,\dd k_2\,\dd k_3\,\dd k_4 ,
\end{equation}
which has the structure of a Boltzmann collision integral: resonant
quartets of waves play the role of colliding pairs, and the kinetic time
scales as $\lambda^{-2}$. The five items are: hierarchy, the equations for
even-order correlations of Fourier amplitudes; limit, weak nonlinearity
with a large box, taken together; retained state, the spectrum; discarded,
phases and higher cumulants; defect, the difference between the measured
spectral transfer and the collision integral. Deng and Hani \cite{DH23}
proved the limit for the cubic nonlinear Schr\"odinger equation over the
full lifespan of the kinetic solution. It is the second canonical track of
the benchmark, chosen because the discarded information is spectral and
not particle-based.

\subsection{Quasilinear theory}

The plasma version closes the interaction between waves and the particle
distribution. Split $f=f_0+\delta f$, with $f_0$ the spatial average, in a
one-dimensional electrostatic plasma. The average of the Vlasov equation is
exact and unclosed,
\begin{equation}
  \partial_tf_0=-\frac{e}{m}\,\partial_v\avg{\delta E\,\delta f}.
\end{equation}
Quasilinear theory \cite{VVS62,DrummondPines62} closes it by using the
\emph{linear} response of $\delta f$ to $\delta E$ about the slowly
evolving $f_0$, which gives a diffusion equation in velocity,
\begin{equation}
  \partial_tf_0=\partial_v\bigl(D(v)\,\partial_vf_0\bigr),
  \qquad
  D(v)\propto\frac{e^2}{m^2}\int|E_k|^2\,\pi\,\delta(\omega_k-kv)\,\dd k,
  \qquad
  \partial_t|E_k|^2=2\gamma_k[f_0]\,|E_k|^2 ,
\end{equation}
with the growth rate $\gamma_k$ given by the slope of $f_0$ at the
resonant velocity. The constant of proportionality depends on the
normalization of the spectrum. The closure assumes many overlapping
resonances and a field that decorrelates before a particle is trapped.

Two connections to earlier sections. The delta function is the same
resonance as in Section~\ref{sec:collisionless}; quasilinear theory is its
nonlinear continuation for the mean distribution. And the natural repair
when the assumptions fail, resonance broadening \cite{Dupree66}, replaces
the delta function by a response function of finite width determined
self-consistently by the diffusion, which is a closure of the
direct-interaction type (Section~\ref{sec:statclosure}). Quasilinear
transport models with fitted saturation rules are in routine use for
turbulent transport in magnetized fusion plasmas, and they are closures in
the sense of Section~\ref{sec:turb}: models, calibrated against nonlinear
simulations.
